\section{课题背景及研究的目的和意义}
\subsection{课题背景}
发展国防工业、微电子工业等尖端技术需要精密和超精密的仪器设备，精密仪器设备要求高速、……

……

\subsection{研究的目的和意义}
气体轴承是利用气膜支撑负荷或减少摩擦的机械构件。……

……

\section{国内外研究现状及分析}
\subsection{国内外研究现状}
国外已经开展了很多关于仿人冗余双臂机器人的研究工作，国内则相对滞后。但在近年来中国科技与经济迅速腾飞的过程中，国内的很多高校及科研机构也在积极开展此方面的研究，具有代表性的有哈尔滨工业大学、清华大学、北京航空航天大学、上海交通大学和北京理工大学等。随着此项技术进一步地向高度的智能化发展，应用领域的不断扩大，机器人正在逐步地将人类从一些危险、复杂、不确定的操作环境中解放出来，改善人类的工作和生活品质，推动社会的进步和发展。

……

\subsection{国内外文献综述及简析}
\subsubsection{冗余机械臂运动学研究现状}
冗余机械臂是指其自身具有的自由度数多于执行任务所需的最少自由度数，实际上具有六个自由度的机械臂，便可以在可达空间内以任意位姿到达任意位置，但是在某些特殊情况下还是会选择可以进行避障、关节力矩优化和奇异回避等优点的七自由度冗余机械臂。

……

\subsubsection{双臂协调轨迹规划研究现状}
早在二十世纪五十年代，最早出现的机器人很多就已经开始采用双臂系统，但广泛应用于工业生产和生活中的，却是便于设计、控制和操作的单臂机器人。

\section{主要研究内容及研究方案}
\subsection{研究内容}
为了确保机器人宇航员能够配合宇航员完成空间站的大部分操作任务，同时保证其能独立完成部分精细操作，机器人宇航员的手臂采用了七自由度冗余结构。而采用冗余结构的机械臂相比于常见的六自由度机械臂（如 PUMA、Stanford），有着许多不可替代的优点，增强了机械臂在执行任务时的可靠性、实用性，可以实现避障、容错和关节力矩优化。但同时也给机械臂的逆运动学求解带来了麻烦，下面将对冗余机械臂的运动学进行详细的分析。

\subsubsection{冗余机械臂建模}
七自由度冗余机械臂的结构形式很多，每种特定的结构也都有自己独特的优势和用途。为了方便机械臂的逆运动学求解，我们采用 S-R-S 形式的七自由度冗余机械臂结构，如图所示。前三个关节交于一点，形成一个球关节（sphere），称作肩关节；后三个关节轴线也相交于一点形成一个球关节，称作腕关节；第四个关节也就是肘关节，对其进行了偏置以便增加肘关节的运动范围，进而增大整个人机系统的工作空间。

……
\addcontentsline{toc}{subsubsection}{……}

\subsection{研究方案}
机器人宇航员的研制包含机械部分、控制系统、视觉反馈以及地面遥操作等一整套系统工程学问题。鉴于该项目涉及的学科广、难度大和任务多，本课题主要负责机器人宇航员七自由度冗余手臂的参数化逆运动学求解，并进行双臂协调操作的轨迹规划，最后搭建参数化的动力学模型进行联合仿真，验证算法的正确性，模拟实现预期的典型操作任务。课题研究流程如下图所示：

……

\section{进度安排及预期目标}
\subsection{进度安排}
本课题的研究时间为 2013 年 9 月开始至 2015 年 7 月份进行答辩，根据之前的项目经验，结合自身学习，科研攻关能力，针对课题研究计划，制定出下面的研究进度：

1. 2013 年 9 月至 2014 年 1 月：开展了机器人方面基础知识学习，了解并掌握了六自由度以及冗余机械臂的逆运动学求解方法。

……

\subsection{预期目标}
预期目标是研制出仿生蛇形机器人完整样机，主要包括机械系统和嵌入式控制系统。本课题预期实现下面四个目标和成果：

1. 完成蛇形机器人的蠕动步态运动，验证规划算法的正确性；

2. 完成蛇形机器人的蜿蜒运动步态，验证规划算法的正确性；

3. 可远程无线控制操作；

……

\section{已具备和所需的条件和经费}
\subsection{实验室条件和经费保障}
目前已经具有一套六自由度机械臂，可以进行轨迹规划算法和力柔顺控制算法的验证。

……

本课题受到载人航天预研项目的大力支持，经费可以得到部分的保障。

……

\subsection{所需条件和经费}
开展今后的课题研究需要一个具有冗余双臂的机器人宇航员作为研究平台，现在正在进行机械部分的采购和加工，其中一支灵巧手已经完成加工和测试。预计年底前完成整机的机械部分组装，随后开展双臂协调操作和力柔顺控制算法的实验验证。

……

\section{预计困难及解决方案}
\subsection{预计困难与技术难点}
课题进行需要学习的知识较多，涉及嵌入式硬件开发、软件开发、电路板设计，机器人的运动学、轨迹规划等知识。

……

\subsection{解决方案}
为了提高课题进展速度和效率，保证学习质量，采用边应用边学习再应用的方法进行主要知识的系统学习，并多动手实际操作，提供动手操作能力，也外后期打下良好的基础。同时采取以下方案：

……

% 学校样例含跳号和省略行；此处手工录入以保持对照文字与编号一致。
% 撰写自己的报告时可替换为 \bibliographystyle{hithesis}\bibliography{reference}。
\begin{thebibliography}{8}
\bibitem{school-lin1992} 林来兴. 空间控制技术［M］. 北京：中国宇航出版社，1992：25-42.
\bibitem{school-xin1999} 辛希孟. 信息技术与信息服务国际研讨会论文集：A 集［C］. 北京：中国科学出版社，1999：45-49.
\bibitem{school-zhao1998} 赵耀东. 新时代的工业工程师［M/OL］. 台北：天下文化出版社，1998［1998-09-26］. \url{http://www.ie.nthu.edu.tw/info/ie.newie.htm}（Big5）.
\item[] ……
\bibitem[7]{school-chen1992} 谌颖. 空间交会控制理论与方法研究［D］. 哈尔滨：哈尔滨工业大学，1992：8-13.
\bibitem[8]{school-kanamori1998} KANAMORI H. Shaking without quaking［J］. Science，1998，279（5359）：2063-2064.
\item[] ……
\end{thebibliography}

% Local Variables:
% TeX-master: "../report"
% TeX-engine: xetex
% End:
